\section{治理建议}
问题一中四大家鱼总量回升时，水草、浮游动物和底栖饵料仍可能下降；问题二中江豚与长江鲟对通道和放流的响应又不相同。治理重点不能只落在捕捞总量，还包括资源层和物种受限环节。监测同时保留四大家鱼总量、底层资源和食压指数，不只报告鱼类总量；鱼道、产卵场连通和放流节律也不按同一方案处理。

问题五的污染情景中，部分过程量短时上升，珍稀物种终态和综合得分却下降。方向相反，所以局部过程改善并不等于系统整体好转。对鄱阳湖等支流湖泊，可将其作为高风险前沿，优先布设入侵物种早期监测、拦截和清除，并把水质、连通性和放流措施纳入同一比较。这里的模型给出措施先后；如果把情景排序直接解释为实际治理效果，结论不在模型能够支持的范围内。

